Model and feature selection used a development partition of 22,056 cases, internally split into 18,747 training
and 3,309 validation cases. The selected configuration was evaluated once on 9,453 temporally later held-out
cases. The 76 routes used in the fixed-replay benchmark belong to this held-out partition and have zero overlap
with model training.

The evaluation dataset is generated by replaying the event log on the BPMN
model. A decision observation is included only when the current event can be
synchronized and fired, the resulting marking enables several possible
successors, and the observed next activity belongs to that candidate set.
This produces $29{,}385$ common held-out decision observations. Detailed split
and replay-coverage statistics are reported in
App.~\ref{app:joao_branching_model_selection}.

\begin{table}[H]
    \caption{Development-selected branching comparison on $29{,}385$ common held-out
    BPMN-replay observations.}
    \label{tab:branching_corrected_methods_joao}
    \centering
    \renewcommand{\arraystretch}{1.05}
    \scriptsize
    \resizebox{\columnwidth}{!}{
    \begin{tabular}{lrrr}
        \toprule
        \textbf{Method}
        & \textbf{Accuracy}
        & \textbf{Macro-F1}
        & \textbf{Weighted-F1} \\
        \midrule
        Majority baseline
        & $0.8385$ & $0.6098$ & $0.8222$ \\
        Random candidate baseline
        & $0.4348$ & $0.3046$ & $0.5365$ \\
        Probabilistic branching
        & $0.7785$ & $0.5147$ & $0.7786$ \\
        Development-selected RF composite
        & $0.9353$ & $0.7073$ & $0.9358$ \\
        \bottomrule
    \end{tabular}
    }
\end{table}

The development-selected RF composite achieves the highest accuracy and macro-F1. Its
macro-F1 of $0.7073$ exceeds the majority-baseline value of $0.6098$, showing
better performance on less frequent branch classes. The probabilistic component
performs below the Random Forest but remains an interpretable Basic method and
the fallback layer of the runtime policy.

The raw classifier, BPMN-constrained policy, and composite runtime produced the
same scores on the common evaluation observations because the highest-ranked
RF prediction was already BPMN-valid in every evaluated row. Fallbacks are
still required during generative simulation, where unseen histories or
different candidate sets may occur.

Results are based on one temporal held-out split; broader generalization is not
claimed.